\documentclass[12pt]{article}

\usepackage[margin=1in]{geometry}
\usepackage{setspace}
\onehalfspacing
\usepackage{natbib}
\usepackage{hyperref}

\title{David Hume's Historical Foundations and James Steuart's Alternative:\\Contrasting Views on Public Debt in Eighteenth-Century Britain}
\author{}
\date{}

\begin{document}

\maketitle

%\doublespacing

\section{Introduction}

The eighteenth century witnessed fundamental debates about the nature and consequences of public debt. Two Scottish thinkers—David Hume and Sir James Steuart—developed sharply contrasting perspectives on government borrowing, shaped by their different readings of history and their exposure to different intellectual traditions. This essay examines how Hume's historical research informed his pessimism about public debt, how his views influenced Adam Smith, and how Steuart arrived at a more positive assessment through engagement with Continental thought and the controversial experiments of John Law.

\section{David Hume's Historical Foundations}

David Hume's concerns about public debt were deeply informed by his research while writing the \textit{History of England} (1754--1762) \citep{hume1754history}. His historical investigation gave him direct exposure to developments that shaped his economic skepticism about government borrowing.

\subsection{The Financial Revolution and Its Consequences}

Hume chronicled England's transformation after 1688 into what historians now call a ``fiscal-military state,'' marked by the founding of the Bank of England in 1694 and the development of a permanent national debt. He observed the dramatic accumulation of debt through successive wars: the War of Spanish Succession, the War of Austrian Succession, and the Seven Years' War. By the time he was writing, Britain's national debt had grown from essentially nothing in 1688 to over £130 million by the 1760s.

\subsection{Specific Historical Episodes}

Several historical events particularly shaped Hume's thinking:

\begin{itemize}
\item The South Sea Bubble (1720) demonstrated the dangers of mixing public debt with speculative finance
\item The growing political power of the ``monied interest''—creditors whose wealth depended on government debt service
\item The increasing burden of debt service on public revenues; by mid-century, debt service consumed roughly half of government revenue
\end{itemize}

\subsection{Hume's Resulting Fears}

In his essay ``Of Public Credit'' (1752, with revisions through the 1760s), Hume predicted that the system would end either in national bankruptcy or in the ``natural death of public credit'' \citep{hume1752public}. He worried that:

\begin{itemize}
\item Debt diverted capital from productive trade and manufacturing to unproductive government consumption
\item It created a rentier class disconnected from productive economic activity
\item It would eventually become unsustainable, leading to either repudiation or ruinous taxation
\end{itemize}

\section{Influence on Adam Smith}

Smith and Hume were intimate friends and intellectual companions in Edinburgh's Scottish Enlightenment circles. The influence flowed through several channels.

\subsection{Direct Intellectual Exchange}

Smith and Hume discussed political economy extensively during the 1750s--1770s. Smith read and admired Hume's essays, including ``Of Public Credit,'' and acknowledged Hume as a major influence on his economic thinking.

\subsection{Smith's Treatment in \textit{Wealth of Nations}}

Book V, Chapter 3 of \textit{Wealth of Nations} (1776) contains Smith's analysis of public debts \citep{smith1776}. In it, Smith:

\begin{itemize}
\item Echoes Hume's historical narrative about the growth of public debt
\item Shares Hume's concern that debt diverts capital from productive investment
\item Worries about the political influence of creditors (``the stock-jobbing and funding interest'')
\item Expresses similar pessimism about the difficulty of ever paying down the debt
\end{itemize}

\subsection{Key Similarity}

Both viewed public debt through the lens of what we might call ``capital displacement''—the idea that government borrowing absorbed capital that would otherwise fund productive commerce and manufacturing. This reflected their transition period's concern with capital accumulation as the driver of economic growth.

\subsection{Historical Validation}

Interestingly, Hume's and Smith's dire predictions didn't materialize—Britain continued to borrow and grew wealthy. Modern economists, including Thomas Sargent, have analyzed this ``puzzle,'' showing how credible commitment to debt repayment, combined with productive investment of borrowed funds and economic growth, allowed Britain to sustain high debt-to-GDP ratios that would have seemed impossible to eighteenth-century observers \citep{sargent_velde1995}.

The Hume-to-Smith transmission represents a fascinating case where historical observation shaped economic theory, even though the theory proved overly pessimistic about institutional capacity to manage public debt.

\section{Sir James Steuart's Contrasting View}

Steuart's \textit{An Inquiry into the Principles of Political Oeconomy} (1767) took a dramatically different position from Hume and Smith \citep{steuart1767}. Rather than viewing public debt as inherently dangerous, Steuart saw it as a potentially productive tool for economic development when properly managed.

\subsection{His Core Argument}

Steuart argued that:

\begin{itemize}
\item Public debt could enhance economic circulation and liquidity
\item Well-managed debt supported both banking and financial market development
\item The key was not whether to borrow, but how to manage the debt and monetary system
\item Public borrowing could ``breathe life into the country's permanent income, by utilizing stagnant funds and trade imbalances'' \citep{stettner1945}
\end{itemize}

\section{Influences on Steuart's Positive View}

\subsection{John Law (1671--1729)—The Critical Influence}

Steuart was profoundly influenced by John Law \citep{law1705}, his fellow Scot who had implemented a radical financial system in France (1716--1720). Unlike Hume and Smith, who viewed Law's Mississippi System as a catastrophic failure proving the dangers of paper money and debt schemes, Steuart had a nuanced ``balanced opinion'' \citep{bentemessek_betancourt2012}.

\subsubsection{What Steuart Learned from Law}

Law's \textit{Money and Trade Considered} (1705) argued that:

\begin{itemize}
\item Money was credit, not just specie
\item Money supply should be determined by the ``needs of trade,'' not gold imports
\item Paper money and credit could stimulate economic activity and employment
\item The problem with Law's System wasn't the concept but the execution—poor management of money, violation of credit rules, and weak banks
\end{itemize}

In his \textit{Principles}, Steuart wrote admiringly:

\begin{quote}
I can point out the utility of banks in no way so striking, as to recal to mind the surprising effects of Mr Law's bank, established in France, at a time when there was neither money or credit in the kingdom. The superior genius of this man produced, in two years time, the most surprising effects imaginable; he revived industry; he established confidence.
\end{quote}

Crucially, Steuart saw Law's attempt to restructure France's massive public debt (over £140 million sterling at Louis XIV's death) through debt-to-equity conversion as conceptually sound, even though it ultimately failed \citep{velde2003, murphy1997}.

\subsection{Charles Davenant (1656--1714)}

Davenant was an English mercantilist writer whose work Steuart explicitly cited \citep{davenant1698}. Davenant had argued:

\begin{itemize}
\item That ``Gold and Silver are indeed the Measures of Trade, but the Spring and Original of it, in all nations is the Natural or Artificial Product of the Country''—wealth wasn't just bullion
\item Favored ``short funds'' (temporary borrowing) over perpetual debt, but accepted borrowing as necessary
\item Saw the development of credit markets as part of commercial progress
\end{itemize}

Steuart references Davenant's views on public borrowing and his proposals for excise taxation as revenue sources. Steuart built on but went beyond Davenant's acceptance of public credit.

\subsection{Continental Cameralist Influences}

During his 20+ years of exile in France and German states (1740s--1763) following the Jacobite rising, Steuart absorbed \textbf{Cameralist} economic thought—the dominant approach in German-speaking Europe. Cameralists like Johann Heinrich Gottlob von Justi emphasized \citep{tribe1993, rossner2016}:

\begin{itemize}
\item Active state management of economic development
\item Credit and public finance as tools of statecraft
\item The state's role in promoting industry and commerce
\item Systematic approaches to public finance
\end{itemize}

This gave Steuart a more positive view of state intervention and public credit than the emerging British classical tradition.

\subsection{French Mercantilist Writers}

Through his time in France, Steuart was exposed to French economic thinking that was more accepting of:

\begin{itemize}
\item State management of finance and trade
\item The productive role of financiers and credit
\item Public debt as part of commercial development
\end{itemize}

\section{Steuart's Synthesis}

Steuart synthesized these influences into a theory where:

\begin{enumerate}
\item \textbf{Debt Creates Liquidity}: Unlike Hume's view that debt diverted capital from productive use, Steuart argued that public debt securities created liquid financial assets that facilitated commerce and banking

\item \textbf{Circulation Matters}: The entire economy benefited from the circulation that debt and paper money enabled—this was the lesson from Law's initial success

\item \textbf{Development Tool}: For ``infant'' trading nations, public credit could help mobilize resources for economic development

\item \textbf{Management Over Prohibition}: The issue wasn't whether to have public debt, but how to manage it properly—maintaining credit rules, sound banking, and honest financial reporting
\end{enumerate}

\section{The Critical Difference}

Where Hume and Smith saw history as warning against debt, Steuart saw it as demonstrating the potential of well-managed public credit systems. The South Sea Bubble and Mississippi System failures weren't proofs that debt was bad—they were lessons in how \textit{not} to manage debt restructuring and monetary policy \citep{menudo2019}.

This made Steuart an outlier in British political economy, closer to what modern scholars call a ``public debt realist'' (along with Alexander Hamilton later) than either a pessimist (Hume/Smith) or an unbounded optimist \citep{salsman2017}. His Continental influences and his sympathy for Law's ambitions positioned him as something of a precursor to later developmental state approaches, where public debt finances productive investment rather than simply displacing private capital.

\section{Conclusion}

The contrast between Hume and Steuart illustrates how the same historical evidence—the Financial Revolution, the South Sea Bubble, Law's System—could be interpreted through different analytical lenses to reach opposite conclusions. Hume, influenced by his narrative of English history and concerned about the rise of a non-productive creditor class, saw public debt as a path to national ruin. Steuart, influenced by Continental cameralism, Law's monetary theory, and mercantilist acceptance of state economic management, saw public debt as a tool for economic development when properly managed.

Both perspectives would echo through subsequent debates about public finance. The Hume-Smith pessimism would dominate British classical economics, while Steuart's more nuanced view would find echoes in Hamilton's American system, German historical school economics, and eventually in modern theories of developmental finance. The puzzle remains that Britain—despite Hume's warnings—managed its growing debt successfully, suggesting that institutional capacity and economic growth can sustain debt levels that pure theory might deem unsustainable.

\bibliographystyle{aer}
\bibliography{hume_steuart_debt}

\end{document}
